\runninghead{The Sunday and Its Texts}
\section{The First Sunday of Advent and Its Texts}
\label{sec:opening}

The Proper of Time in the Roman Missal opens on its first page with a man
lifting his soul to God:
\latin{Ad te levávi ánimam meam}, ``To thee have I lifted up my soul.'' The
words open Psalm 24 (25), and they open the Church's year. A few minutes
later the Gospel describes the end of the world. In St Luke's Gospel Christ tells
his disciples of signs in the sun, the moon and the stars, of nations in
distress at the roaring of the sea, of men ``withering away for fear and
expectation of what shall come upon the whole world'' (Lk 21:25--26). Then he
gives a command that seems to belong to another scene altogether: ``when
these things begin to come to pass, look up and lift up your heads, because
your redemption is at hand'' (v.~28). The heads lifted in the Gospel and the
soul lifted in the Introit are lifted by the same Latin verb, \latin{leváre}.

The Missal heads the Mass with its Roman station, \latin{Statio ad S.
Mariam maiorem}, St Mary Major. Blessed Ildefonso Schuster
recalls that this basilica has venerated, since Pope Sixtus III in the fifth
century, a Roman reproduction of the crib of Bethlehem, and he takes the
station to point the faithful toward the manger at the start of the season
(\work{The Sacramentary}, vol.~1, p.~321). The Gospel read there is of the Son
of man ``coming in a cloud, with great power and majesty.''

Between the psalm and the Gospel St Paul tells the Romans that ``it is now
the hour for us to rise from sleep. For now our salvation is nearer than when
we believed'' (Rom 13:11). The Collect, addressed to the Son, asks him to
stir up his power and come. Two further chants come from Psalm 84 (85): the
Alleluia asks God to show his mercy and give his salvation, and the
Communion promises that ``the Lord will give goodness: and our earth shall
yield her fruit.'' The Secret asks that the sacrament cleanse us so that we
may come purer to its source, and the Postcommunion that the mercy received
in the temple may prepare us for the coming feast of our restoration.

The texts raise questions. How can signs that make men wither be the moment
for the disciples to lift their heads? Are the sun, the moon and the stars of
the Gospel the sky above us, or figures of something nearer? What is the
salvation that is ``nearer than when we believed,'' and how does a man who
hears that it is the hour to rise get up, if he would rather sleep? What do
the Psalm 84 chants ask for when they ask God to ``show us thy mercy,'' and
what is the fruit our earth is promised?

Four interpretations of the whole Mass follow, each built on the reasoning of
Fathers and saints who expounded one of its texts. The first follows St
Gregory the Great's homily on this Gospel, in which the command to lift the
head means to gladden the heart because a world we are not to love is
ending, and St Cyril of Alexandria's teaching that the redemption is the
body's resurrection and creation made new. The second follows St Augustine
and St Ambrose, who read the failing lights of heaven in the Church under
persecution and find Christ's coming already under way in her, Augustine
in his body the Church and Ambrose, with a ``perhaps,'' in each believer who
receives him; for Ambrose the sun is Christ and the moon the Church, for
Augustine the Church is sun, moon and stars at once. The third begins from the
Epistle with St John Chrysostom, who compares Paul's readers to runners near
the goal, and St Thomas Aquinas, who distinguishes the ways
salvation can be near and adds how the Church hears the Apostle's words in
Advent. The fourth begins from Psalm 84 with Augustine, for whom God has given
us Christ as man and will give the same Christ as God, and St Robert
Bellarmine, for whom the psalm asks for the Saviour's coming and the earth
that yields its fruit is the Virgin. The first two interpretations differ on
what the Gospel's signs are; the last two start from other texts and meet the
Gospel at other points. The Scriptural
Date and Location appendix (\ref{sec:date-location}) gives the setting and
date of each passage the Mass reads or sings.

\begin{maptable}
Introit &
\latin{Ad te levávi}; Ps 24:1--3 (25:1--3); verse Ps 24:4 (25:4) &
The psalmist lifts his soul to God, trusts not to be put to shame, and
says that none who wait for God will be confounded. The verse asks God to
show his ways and teach his paths.\\
Collect &
\latin{Excita, quǽsumus, Dómine, poténtiam tuam, et veni}; addressed to the
Son &
The Church asks Christ to stir up his power and come, so that by his
protection and deliverance we may be rescued and saved from the dangers that
threaten us because of our sins.\\
Epistle &
\latin{Fratres: Sciéntes, quia hora est}; Rom 13:11--14a &
Paul tells the Romans it is the hour to rise from sleep, that the night is
past and the day at hand, and that they are to cast off the works of
darkness and put on the Lord Jesus Christ.\\
Gradual &
\latin{Univérsi, qui te exspéctant}; Ps 24:3--4 (25:3--4) &
None who wait for God will be confounded; the verse asks God to make his ways
known.\\
Alleluia &
\latin{Osténde nobis, Dómine}; Ps 84:8 (85:8) &
``Shew us, O Lord, thy mercy; and grant us thy salvation.''\\
Gospel &
\latin{Erunt signa in sole}; Lk 21:25--33 &
Christ foretells signs in the heavens and distress on earth, his coming in a
cloud with power and majesty, the command to lift up the head, the fig tree,
the generation that will not pass and the words that will not pass.\\
Offertory &
\latin{Ad te levávi}; Ps 24:1--3 (25:1--3) &
The Introit's antiphon sung again at the offering.\\
Secret &
\latin{Hæc sacra nos, Dómine} &
May these holy things, cleansing us by their power, make us come purer to
their source.\\
Communion &
\latin{Dóminus dabit benignitátem}; Ps 84:13 (85:13) &
The Lord will give goodness, and our earth shall yield her fruit.\\
Postcommunion &
\latin{Suscipiámus, Dómine, misericórdiam tuam} &
May we receive God's mercy in the midst of his temple, and go before the
coming feast of our restoration with fitting honour.\\
\end{maptable}
